\documentclass[12pt]{jlreq}
\usepackage{amsmath, amssymb}

\newcommand{\C}{\mathbb{C}}
\newcommand{\R}{\mathbb{R}}
\newcommand{\Z}{\mathbb{Z}}
\newcommand{\Tr}{\operatorname{Tr}}
\newcommand{\Hom}{\operatorname{Hom}}
\newcommand{\GL}{\operatorname{GL}}
\newcommand{\rank}{\operatorname{rank}}
\newcommand{\Id}{\operatorname{Id}}
\newcommand{\Ker}{\operatorname{Ker}}

\begin{document}

$D$ は $\R^n$ 内の有界閉集合，$\varphi$ は $\R$ 上の下半連続な非負値関数とし，ルベーグ可測関数 $f : D \to \R$ に対して
\[
  \Phi(f) = \int_D \varphi(f(x)) \, dx
\]
とおく。$\Phi(f)$ は $+\infty$ の値をとることも許すものとする。このとき，次を示せ。

\begin{enumerate}
  \item $\Phi$ は $C(D)$ 上で下半連続である。
  \item $\Phi$ は $L^p(D)$ 上で下半連続である。ただし $1 \le p < \infty$ とする。
\end{enumerate}
ここで $C(D)$, $L^p(D)$ は，それぞれノルム
\[
  \|f\|_\infty = \sup_{x \in D} |f(x)|,\quad \|f\|_p = \left( \int_D |f(x)|^p \, dx \right)^{\frac{1}{p}}
\]
から決まる位相をもつ $D$ 上の関数空間である。

\end{document}